% !TeX root = ../papers/nc-p3.tex
\section{二つの行列式二次式が消えるファイバー}
\label{sec:klein-zero-quadratics}
節\ref{sec:klein-elliptic-completion}の薄い成分が非零で、$z\ne\pm1$、
交差ブロックと組内のSchur主ブロックが可逆なチャートを用いる。
特に$z\lambda(1-\lambda^2z)D_0\ne0$である。

\begin{theorem}\label{thm:klein-all-pairings}
このチャートで二つの隣接階数が6となるテンソルは、五次元楕円曲線族の閉包と
定理\ref{thm:klein-other-pairings}の四つの表示で全て覆われる。
$P,Q$がともに零多項式のとき、共通倍率を$a=1$で固定すると八種類の基底条件があり、
それぞれのファイバーは単一の双一次方程式で決まる。
\end{theorem}
\begin{proof}
係数環は有理数上の多項式環をチャートの分母で局所化したものである。
$P(x)$と$Q(y)$はそれぞれ原始である。実際、定数項は順に
\[
 n=(a-zd)(za-d)/(z+1)^2,\qquad o=(b-c)(z^2b-c)/(z+1)^2
\]
で、各一次因子を零にしても対応する最高次係数は恒等的に零でない。
例えば$(a,b,c,d,z,\lambda)=(1,3,4,5,2,3)$から順に$d=za,a/z$、
$c=b,z^2b$と一箇所を置換すればよい。共通判別式は分数体で平方でない。
従ってGaussの補題により、$P(x)=0$と$Q(y)=0$の各全空間は整である。
両者の一般点から楕円曲線族への有理逆写像は
$2px+m=-(2qy+k)$で与えられる。

$\delta=0$で片方だけが零多項式の場合を考える。例えば$P\equiv0$なら、
解を持つ非零の$Q$は$q\ne0$の二重根を持つ。このとき
$y=-(2px+m+k)/(2q)$は$P(x)=0$上で正則であるため、任意の$x$に
付随する解は一般の族の特殊化である。逆側でも同様である。

$P\equiv Q\equiv0$なら上の定数項から$d\in\{za,a/z\}$,
$c\in\{b,z^2b\}$となる。$a=0$または$b=0$では$m=1$または$k=-1$となるので、
$a,b$は非零である。$a=1$とする。$f_0 g_0 e_0\ne0$であり、
$p=q=0$と$m=k=0$から、$f_1,g_1$の一方だけと$e_1,h_1$の一方だけが零である。
表の各対はこの零となる係数で、隣は$b$の値である。
\[
\begin{array}{cc|cc|cc}
 d&c& \multicolumn{2}{c|}{\text{第一の場合}}&\multicolumn{2}{c}{\text{第二の場合}}\\\hline
 z& b &(f_1,e_1)&1/\lambda &(g_1,h_1)&\lambda z\\
 z& z^2 b &(f_1,h_1)&\lambda &(g_1,e_1)&1/(\lambda z)\\
 1/z& b &(f_1,h_1)&1/(\lambda z)&(g_1,e_1)&\lambda\\
 1/z& z^2 b &(f_1,e_1)&\lambda/z &(g_1,h_1)&1/(\lambda z^2)
\end{array}
\]
検算の入力を明記する。$f_1=0$なら$T_{00}=T_{11}=1/f_0$、
$g_1=0$なら$(T_{00},T_{11})=(1,z^2)/g_0$である。
$e_1=0$なら$(T_{01},T_{10})=(1,z)/e_0$、
$h_1=0$ならこれが$(z,1)/e_0$となる。これらと
$T=Y_\lambda A_0^{-1}X_\lambda$を等置し、上記の非零因子で飽和させる。
十六の選択のうち八つでは単位イデアルを得る。他の八つでは表示した$b$を
強制する式がイデアルに入り、代入で全式が消える。
\path{computations/verify_klein_zero_quadratics.py}が有理係数のGröbner基底で検算する。

各行の代入後、$e,h,f,g$を前の一次式で表し、$F_\lambda(A_1)$に
要求される四式の分母を払う。零でない式は全て、係数環の単元倍を除いて同じ
$H(x,y)=\alpha xy+\beta x+\gamma y+\eta$であり、$\alpha$は単元である。
この主張も同コードが各式の係数を照合し、倍率の因子が全て可逆であることを検査する。
例えば表の左上の場合、単元倍を選ぶと
\[
 H=\lambda^2(\lambda^2 z-1)xy
   -\lambda^2(z-1)^2(x+y)+(z-1)^2(\lambda^2z-1)
\]
となる。ここでも$x,y,\det A_1$の非零条件を保つ。

一般に残る四式は$y$について線型である。その一つの$y$係数が基点で非零なら、
整な$P(x)=0$上の近傍でその式から$y$を解ける。一般点でその$y$は楕円曲線族の
ものであり、残る全ての等式は整性から近傍に延びる。よって当該点はその族の閉包に入る。
$x$係数が非零なら$Q(y)=0$から同様に従う。各固定ファイバーの$H$は$xy$係数が
非零なので既約な双一次曲線か、相異なる縦横二直線である。各成分の稠密開部では
少なくとも一方向に解け、交点も閉包に入る。この議論と
定理\ref{thm:klein-other-pairings}でチャート全体の被覆を得る。
ここで決めたのは階数軌跡であり、閉包上の各閉点の全次数正則性はなお別の判定を要する。
\end{proof}
